% Nature Biotechnology Brief Communication: neutral source-inspired sample.
% Based on a supplied edition published online in 2025 and in a 2026 issue.
% All prose, equations, diagram labels and metadata below are demonstration content.
\PassOptionsToPackage{no-math}{fontspec}
\documentclass[UTF8,fontset=none,twocolumn,scheme=plain]{ctexart}
\usepackage[paperwidth=595.276bp,paperheight=790.866bp,left=39.685bp,right=35.716bp,top=58bp,bottom=43bp,headheight=16bp,headsep=9bp,footskip=24bp]{geometry}
% Shared dependencies and neutral material. Compile with XeLaTeX.
% No publisher logos, paper text, data, figures, or download metadata are included.
\usepackage{amsmath,amssymb,graphicx,xcolor,tikz,fancyhdr,needspace,xurl,hyperref,balance,etoolbox}
% The legacy balance package omits the first split box's final baseline depth.
% Reserve that measured depth rather than hiding a CJK overfull-box warning.
\makeatletter
\patchcmd{\@BAlancecol}{\setbox\@leftcolumn=\vbox to \@tempdima}{\advance\@tempdima by \dp1\relax\setbox\@leftcolumn=\vbox to \@tempdima}{}{\PackageWarning{cns-template}{Balance depth patch unavailable; inspect final columns}}
\makeatother
\usetikzlibrary{arrows.meta,positioning}
\usepackage[protrusion=true,expansion=false]{microtype}
\IfFontExistsTF{Noto Serif CJK SC}{%
  \setCJKmainfont{Noto Serif CJK SC}%
}{%
  \setCJKmainfont[BoldFont=FandolSong-Bold.otf,ItalicFont=FandolKai-Regular.otf]{FandolSong-Regular.otf}%
}
\IfFontExistsTF{Noto Sans CJK SC}{%
  \setCJKsansfont{Noto Sans CJK SC}%
}{%
  \setCJKsansfont[BoldFont=FandolHei-Bold.otf]{FandolHei-Regular.otf}%
}
\setmainfont{TeX Gyre Termes}\setsansfont{TeX Gyre Heros}
\setmonofont{TeX Gyre Cursor}
\definecolor{JournalBlue}{HTML}{0097C9}
\definecolor{JournalRed}{HTML}{B83443}
\definecolor{ScienceRed}{HTML}{CD1021}
\definecolor{CiteBlue}{HTML}{0000EE}
\hypersetup{colorlinks=true,linkcolor=JournalBlue,urlcolor=JournalBlue,
  pdftitle={CNS论文汉化：中性版式示例},pdfauthor={示例模板；非期刊官方模板},
  pdfsubject={Neutral source-inspired Chinese layout demonstration; no real research claims}}
\urlstyle{same}\raggedbottom
\setlength{\parindent}{1em}\setlength{\parskip}{0pt}
\setlength{\emergencystretch}{2em}\tolerance=2200
\XeTeXlinebreaklocale "zh"\XeTeXlinebreakskip=0pt plus .1pt
\pagestyle{fancy}\fancyhf{}
\renewcommand{\headrulewidth}{0pt}\renewcommand{\footrulewidth}{0pt}
\newcommand{\PaperTitle}{中文题名占位：从组织结构到证据解释}
\newcommand{\SampleAuthors}{Author A\textsuperscript{1}, Author B\textsuperscript{2,*}}
\newcommand{\SampleAffiliations}{\textsuperscript{1}作者单位一（示例占位）。\quad
  \textsuperscript{2}作者单位二（示例占位）。\par
  *通信作者信息待填。DOI、收稿和出版日期均应从目标原文核实。}
\newcommand{\SampleAbstract}{本段是中性摘要占位文字，不陈述任何真实研究发现。它演示中文题名、作者信息、摘要和双栏正文之间的层级关系。实际使用时，应完整替换为目标论文的忠实中文译文，并保留研究对象、方法、结果和结论的限定条件。图注、参考文献、作者声明和来源信息均属于独立的覆盖项。模板中的图形仅说明排版结构，不表示生物学机制、实验结果或统计关系。}
\newcommand{\DemoCaption}{\textbf{图1. 中性证据链示意（占位图，无实验数据）。}
  \textbf{A}，输入节点表示需要核对的原文材料。
  \textbf{B}，处理节点表示逐段转换及复核。
  \textbf{C}，输出节点表示正文、图注和书目等独立交付部分。
  箭头仅表示本示例的阅读顺序，不代表因果关系、效应方向或数量变化。
  该图由模板自行绘制，所有节点均为文字占位；无样本量、误差线、统计检验或比例尺。
  实际论文图注应逐项保留原稿报告的样本量、单位、检验方法、限定及图形署名。}
\newcommand{\DemoVector}{%
 \resizebox{\linewidth}{!}{%
 \begin{tikzpicture}[font=\sffamily\small,node distance=18mm,
   box/.style={draw=black!70,rounded corners=2pt,fill=black!3,minimum width=26mm,minimum height=16mm,align=center},
   flow/.style={-{Stealth[length=2.5mm]},line width=.7pt}]
  \node[box] (a) {\textbf{A}\quad 输入\\原文材料};
  \node[box,right=of a] (b) {\textbf{B}\quad 处理\\转换与复核};
  \node[box,right=of b] (c) {\textbf{C}\quad 输出\\中文译稿};
  \draw[flow] (a)--(b);\draw[flow] (b)--(c);
  \node[below=11mm of a,align=center] {标题与摘要\\正文与图表};
  \node[below=11mm of b,align=center] {保留限定\\检查单位};
  \node[below=11mm of c,align=center] {可追溯引文\\明确原文信息};
  \node[below=29mm of b,font=\sffamily\footnotesize,text=black!65] {示意占位；不包含真实实验数据};
 \end{tikzpicture}}%
}
\newcommand{\SampleBibliography}{%
 \section*{参考文献（示例占位）}
 {\fontsize{7.6bp}{10bp}\selectfont\setlength{\parindent}{0pt}\interlinepenalty=10000
 \hypertarget{ref:1}{}1. Author A. (20XX). Example reference title. \textit{Journal placeholder}. 卷、页码与DOI待填。此条仅供演示引文跳转。\par\smallskip
 \hypertarget{ref:2}{}2. Author B. (20YY). Another example reference title. \textit{Journal placeholder}. 出版信息待填；不要把占位条目带入正式译稿。\par}
}
\newcommand{\SampleSourceInfo}{%
 \section*{原文信息}
 以下字段全部为示例占位：作者；英文原标题；期刊；年份；卷期；页码或文章号；DOI或来源链接。实际使用时必须逐项替换并核对。\par
 本文件是\JournalName 风格的非官方中文译稿样板，依据一次已核验文章版式进行结构化改编，不代表该期刊所有年代或文体的统一标准，也不保证像素级复刻。无真实论文内容或期刊标识图像。\par
}
% Deliberately neutral demo prose, not translated from a published paper.
\newcommand{\SampleOpening}{%
 \section*{引言：示例正文与引文}
 这是一段用于检查双栏阅读节奏的中性示例文字。段落中的术语、缩写和标点需要在实际项目中与源文保持一致。中文译稿应呈现原论文的完整论述，而不是把多段证据压缩成一段摘要。这里的引文只链接到本文件末尾明确标为占位的参考文献\SampleCiteOne。\par
 若原句仅表示相关、可能或提示，译文也应保留相同的证据强度。数字、单位、正负号、上下标和范围端点应逐项核对。原文存在内部矛盾时，应分别保留相应位置的陈述，再用简短译注说明；不能根据常识把不同的数字静默统一。\par
 \subsection*{内容层级与跨页衔接}
 此小标题演示二级层级。实际论文可能采用不同的小标题密度，模板不会替代对原文结构的识别。双栏文章需按左栏、右栏和跨页续句的顺序读取。若段落在页末中断，应在完成整句核对后接入下一页，避免把图中文字误当正文，或重复翻译同一段。\par
 图表具有独立的信息范围。正文引用到某幅图，并不意味着其图注可以省略。图注中的每个分面、比较对象、时间点、样本量及统计说明都应对应源文。英文图内标签可以保留，并在图旁或文末提供中文释义；实际采用哪种安排，需要依据图形尺寸和译文长度判断\SampleCiteTwo。\par
 \subsection*{来源与可维护数据}
 本段继续演示正文密度。可编辑文本和图注数据应与最终版面分开维护，以便修订时准确定位变化。若只调整版式，应能证明正文、图注和书目内容没有被意外改写。更换字体或纸张后，须重新检查断行、图注与图片的相对位置以及页脚可读性。\par
 文内引用与书目之间应有明确对应关系。这里的两条演示书目没有真实DOI，也没有虚构的实验信息。使用模板时可增删条目，但必须同时更新编号、链接目标和引用位置。页码只表示中文译稿自己的分页，不用于冒充原出版页码。\par
 \subsection*{结果叙述的演示占位}
 本节的文字仅用于呈现长段落在双栏页面中的流动，不报告真实结果。实际翻译应保留观察对象、实验条件和结局指标之间的关系。例如，原文分别描述体外实验与动物实验时，不能把两类结果合并为同一层面的结论；研究使用的对照、时间顺序和分析范围也必须保留。图表中的比较范围不应扩大为对整个人群或全部疾病阶段的判断。\par
 论述中的转折与限制同样属于正文。若研究先提出一种解释，再讨论另一种可能，译文应沿用这一推进顺序。作者明确写出的不确定性应能在中文中被识别；如果尚未确定某个现象的方向或机制，就不能改成已经建立的因果关系。对于术语近似但研究层级不同的概念，可在首次出现时保留缩写，后文保持统一。\par
 \subsection*{方法信息与可复核性}
 方法部分往往包含密集数字和专有名词。翻译时须区分浓度、总量、频率和处理时长，不能因为中文句式调整而改变参数的归属。试剂名称、软件版本和数据访问标识通常保留原样；上下限、排除规则和显著性标准则需要逐字符核对。这里没有提供任何真实实验参数，避免把模板占位误认为可以直接执行的方法。\par
 表格与补充材料的引用也需要保留原层级。正文对补充图的指向并不意味着模板中已经附有该补充图；若源文件只包含主文，应如实标注所交付的范围。图表编号、附录顺序和参考文献数量都应由目标文章决定，本示例中的结构不构成固定配额，也不要求所有论文采用同样的章节。\par
}
\newcommand{\SampleClosing}{%
 \section*{讨论：版式适配与范围}
 上方矢量图是模板原创的排版占位图，不来自任何论文。图形与完整演示图注被放在同一页，便于读者直接核对分面说明。真实项目应使用经过授权或用户提供的原图素材，并保留图注署名与许可信息。尺寸不足时，可适度缩放或调整栏宽，不以删除图注解决空间问题。\par
 模板仅记录一次已核验文章的可观察设计特征。即使文章出自同一期刊，不同年代、研究论文与综述、校正清样与正式版本也可能采用不同的页眉和引文形式。应以用户提供的实际PDF为准，不能仅凭期刊名称推断配色、字号或引用方式\SampleCiteOne。\par
 \subsection*{演示检查项}
 本样板用于展示可编译性、中文字形、标题层级、跨栏图注、页眉页脚和内部链接。图中没有实际数据，因此不应被当作科学插图引用。替换内容后，还需检查所有页的渲染结果，确认无截字、碰撞、异常空白和遗漏段落。\par
 \SampleBibliography
 \SampleSourceInfo
}


\usepackage{array,booktabs,tabularx,eso-pic}
\definecolor{BiotechYellow}{HTML}{E4DC25}
\definecolor{BiotechBlue}{HTML}{0068A6}
\newcommand{\JournalName}{Nature Biotechnology}
\hypersetup{linkcolor=BiotechBlue,urlcolor=BiotechBlue,
  pdftitle={Nature Biotechnology：简讯中性中文样板}}
\setlength{\columnsep}{12bp}
\renewcommand{\normalsize}{\fontsize{9bp}{12.1bp}\selectfont}
\ctexset{section={format=\sffamily\bfseries\fontsize{10.5bp}{13bp}\selectfont,beforeskip=9bp,afterskip=3bp},subsection={format=\sffamily\bfseries,beforeskip=7bp,afterskip=2bp}}
\newcommand{\SampleCiteOne}{\textsuperscript{\hyperlink{ref:1}{\color{BiotechBlue}1}}}
\newcommand{\SampleCiteTwo}{\textsuperscript{\hyperlink{ref:2}{\color{BiotechBlue}2}}}
\fancyhead[L]{\sffamily\fontsize{9bp}{11bp}\selectfont 简讯\quad Brief Communication}
\fancyhead[R]{\sffamily\fontsize{7.5bp}{10bp}\selectfont\textcolor{BiotechBlue}{DOI：待原文核实}}
\renewcommand{\headrulewidth}{.3bp}
\renewcommand{\footrulewidth}{.3bp}
\fancyfoot[L]{\sffamily\fontsize{7bp}{9bp}\selectfont\textcolor{BiotechBlue}{Nature Biotechnology}\quad |\quad 非官方中文样板}
\fancyfoot[R]{\sffamily\bfseries\fontsize{7bp}{9bp}\selectfont 中文译页\ \thepage}
\fancypagestyle{opening}{\fancyhead{}\renewcommand{\headrulewidth}{0pt}}
\newcommand{\DemoParaA}{这段原创中性文字展示简讯正文的排版密度，不报告真实研究结果。实际译稿从原文首段开始，原稿没有标题时不增加“引言”标题。读者应能区分研究目的、评价对象与适用范围；译者需要保留作者限定的比较对象，并让段落顺序与原文一致。这里的数字上标仅链接本文件的演示书目，不对应任何已发表研究\SampleCiteOne。\par}
\newcommand{\DemoParaB}{不同输出之间的比较，需要保留各自的定义与评价范围。原文若将样本、分组、标签和表示形式分别定义，中文也应逐一对应，不因词语相近而合并。某项指标的方向取决于其具体定义，分数较高是否代表更好不能从名称猜测。引用图表时，应保留原稿的面板字母及对应关系；不能用一句结论替代完整的方法说明。\par}
\newcommand{\DemoParaC}{示例公式中的符号只用于检查数学排版。设示例值为 $q_i$、非负权重为 $w_i$，且权重之和非零，下面的表达式演示分式、求和、上下标与公式编号；它不是目标论文的公式，也不表示经过验证的评价方法。实际使用时应将公式、符号定义及编号一并替换，并检查正文中的引用位置。\par}
\newcommand{\DemoParaD}{逐栏阅读时，页末的句子可能在下一页主图之后才继续。恢复这一顺序需要同时核对段落缩进、图注边界及后续句子，不能把图中坐标文字接入正文。对图内英文的处理遵从本次任务要求；保留英文时，应在适当位置给出完整中文释义。正文、图注和附录的覆盖清单应分别维护，避免其中一项已完成而其他部分被遗漏。\par}
\newcommand{\DemoParaE}{研究对象和适用范围的限定同样属于正文。只有在原文明确支持时，才可以把观察结果表述为因果关系，或把特定条件下的结果推广至其他场景。转折句、例外情况、无统计学显著性的发现及作者提出的限制都应完整译出。示例不提供实验数值，真实项目中的单位、比例、置信区间和软件版本须逐项核对。\par}
\newcommand{\DemoParaF}{简讯的篇幅较短，并不意味着方法、扩展数据、作者声明或出版许可可以省略。原文件若把这些材料排在连续页码之后，仍需识别它们属于当前 PDF 的范围。只被正文提及而没有提供的补充文件，应明确区分。这里的人工分页仅用于三页样板展示，正式文章应按照中文长度重新安排\SampleCiteTwo。\par}
\newcommand{\DemoParaG}{来源版本还可能影响正文与附录的组织方式。在线发表版本、正式卷期版本及作者接受稿可以具有相同题名，却采用不同页码和排版。处理时应记录当前实际提供的文件版本，分别核对在线日期、卷期和版权年份。某个字段没有出现在源文中时，可以明确保留待核实状态；不应根据文件名或相邻文章补出未经证实的信息。\par}
\newcommand{\DemoParaH}{正文中的比较关系往往需要跨越多个句子才能完整理解。前一句可能先定义评价条件，后一句才给出结论；也可能在随后一段补充反例或使用限制。为此，术语表只统一用词，不代替逐段判断。保留段落、转折和引文位置，有助于读者回到原文核查。中性样板只演示这些信息的承载方式，不能证明任何具体译稿已经准确。\par}
\newcommand{\DemoParaI}{图形与版式之间存在实际的尺寸约束。横向长图、纵向流程图和多面板数据图可以采用不同安排，但每种安排都应保留标签的阅读顺序、刻度和完整图注。必要的文字释义可单独分组，组名与面板字母需清楚对应。缩放时首先观察最小标签能否辨认，再调整图页或图注布局；仅使外框落入栏宽并不足以完成核验。\par}
\begin{document}\normalsize
\thispagestyle{opening}
\AddToShipoutPictureBG*{\AtPageUpperLeft{\put(0,0){\begin{tikzpicture}[remember picture,overlay]
\draw[BiotechYellow,line width=3bp] ([xshift=39.685bp,yshift=-36bp]current page.north west)--([xshift=-35.716bp,yshift=-36bp]current page.north east);
\end{tikzpicture}}}}
\twocolumn[{%
\noindent{\sffamily\fontsize{16bp}{20bp}\selectfont Nature Biotechnology\par}\vspace{22bp}
\noindent{\sffamily\bfseries\fontsize{9bp}{11bp}\selectfont 简讯\quad Brief Communication\hfill\mdseries\fontsize{7.5bp}{10bp}\selectfont\textcolor{BiotechBlue}{DOI：待原文核实}}\par
\vspace{4bp}\hrule height .3bp\vspace{7bp}
\noindent{\bfseries\fontsize{25bp}{31bp}\selectfont 中文简讯题名占位：评价方法、\par 证据边界与可复核说明\par}\vspace{24bp}
\hrule height .3bp\vspace{5bp}
\noindent\begin{minipage}[t]{.32\textwidth}\vspace{0pt}\sffamily\fontsize{8bp}{11bp}\selectfont
收稿日期：待核实\par\vspace{5bp}\hrule height .3bp\vspace{5bp}
接受日期：待核实\par\vspace{5bp}\hrule height .3bp\vspace{5bp}
在线发表：待核实\par\vspace{5bp}\hrule height .3bp\vspace{5bp}
非官方中文译稿样板\par
\end{minipage}\hfill
\begin{minipage}[t]{.656\textwidth}\vspace{0pt}
{\sffamily\bfseries\fontsize{9bp}{12bp}\selectfont\SampleAuthors\par}\vspace{6bp}\hrule height .3bp\vspace{6bp}
{\fontsize{10.5bp}{14.4bp}\selectfont\SampleAbstract\par}
\end{minipage}\par\vspace{8bp}
{\sffamily\fontsize{7.1bp}{9.4bp}\selectfont\SampleAffiliations\par}\vspace{12bp}
}]
\balance
\DemoParaA\DemoParaB\DemoParaC
\begin{equation}Q=\frac{\sum_{i=1}^{m}w_iq_i}{\sum_{i=1}^{m}w_i}\label{eq:demo}\end{equation}
式（\ref{eq:demo}）中的所有字母均为排版占位变量。引文上标与公式编号承担不同功能，不应相互替换。若原文另有符号约定，应保留原约定及适用条件。\par
\DemoParaD\DemoParaE\DemoParaG\DemoParaH
\clearpage
\twocolumn[{%
\begin{center}\begin{minipage}{.84\textwidth}\DemoVector\end{minipage}\end{center}
\vspace{3bp}\noindent{\sffamily\fontsize{7.7bp}{10.3bp}\selectfont\DemoCaption\par}\vspace{10bp}\hrule height .3bp\vspace{11bp}
}]
\balance
\DemoParaD\DemoParaB\DemoParaE\DemoParaF
图注和图像作为一个完整信息单元在本页呈现。实际图注若超过一页可容纳的高度，应改用清楚的延续标记或专门图页，并核对所有面板说明。模板中的图没有任何科学观测值；其色块、箭头和节点仅帮助检查宽度、对齐与可读性，不应被解释为真实数据。\par
数学符号可与中文段落共同排版，但不能根据图示外观重建不存在的数值。这里再用一个完全中性的恒等关系演示第二个编号：\par
\begin{equation}\bar q=\frac{1}{m}\sum_{i=1}^{m}q_i,\qquad m\ge1.\end{equation}
本式仅用于测试字号、基线和编号对齐，不能替代实际研究方法。若公式包含分段条件、向量或矩阵，应依据原稿扩展环境，并重新检查双栏边界。较长公式可以合理断行，但不可删除项、改变括号作用范围或移走条件。\par
\DemoParaI\DemoParaH
\section*{扩展材料的衔接示例}
此处演示从主文转入扩展材料的逻辑连接。扩展图、报告摘要与数据声明是否存在，取决于实际提供的源文件；模板不会自动生成这些内容。附录中的选择框、表格字段和自由文本答案都属于独立信息，不能只保留表题或栏目名。\par
\clearpage
\twocolumn[{%
\section*{扩展材料与报告表（中性示例）}
\noindent 此表只演示报告类附录的字段和值，不填写真实研究结论。正式译稿应逐项保留源表字段、状态、说明及未适用项。\par\vspace{8bp}
\begingroup\sffamily\fontsize{8bp}{11bp}\selectfont
\setlength{\tabcolsep}{6bp}\renewcommand{\arraystretch}{1.3}
\noindent\begin{tabularx}{\textwidth}{@{}p{89bp}p{87bp}X@{}}
\toprule\textbf{字段（示例）}&\textbf{状态占位}&\textbf{原文说明占位}\\\midrule
资料范围&待原文核实&列出本次实际提供的主文、扩展图、表格及报告摘要，不推定另有文件。\\
样本与分组&待原文核实&保留样本定义、分组依据和排除说明；没有原始数值时不补写。\\
统计说明&待原文核实&分别记录检验名称、比较对象及适用条件，不使用模板默认阈值。\\
数据可用性&待原文核实&填写源文提供的访问说明与标识；示例不提供真实数据链接。\\
代码可用性&待原文核实&填写版本、获取方式和限制；代码未提供时不宣称已经复现。\\
报告字段&$\square$ 状态待核实&空方框只演示符号字形，不代表作者作出的选择；使用时按原稿替换。\\
材料获取&待原文核实&保留源文给出的获取条件、通信对象和使用限制，不用模板许可替代。\\
分析边界&待原文核实&区分预先设定的分析与探索性分析，原文未说明时不替作者推定。\\
作者声明&待原文核实&分别对应贡献、资助和利益冲突字段，不合并不同主体的声明。\\
附录定位&待原文核实&保留原附录编号与主文引用关系；未提供的文件不能列作已译内容。\\
\bottomrule
\end{tabularx}\par\endgroup\vspace{14bp}
}]
\balance
\section*{方法与声明占位}
附录中的方法并非主文结果的摘要。操作步骤、参数设置、检验假设和数据访问条件都应在正式译稿中完整呈现。若原文明确说明某项分析未进行，不能因为报告表留有该字段而补出一项分析。\par
报告字段的中文长度可能超过原文，因此本表允许单元格自然换行。字段及其回答应作为一个整体保留，不能因原版留白较窄而截去说明。若完整表格跨页，应为续页重复栏目名，并在末尾保留缩写和脚注。此页仅展示短表，不预设真实论文需要多少行或多少个选项。\par
原作者、出版者和图像权利信息应按目标原文保留；中性模板的许可不改变用户另行提供材料的权利归属。本文档仅演示排版，不授予对任何原刊 logo、图像或文字的额外使用权。\par
\section*{版式与版本说明}
本样板提炼一篇简讯的观察结果：首页使用黄色顶线、左日期窄栏与右作者摘要宽栏，普通页面使用细横线与稿件类型/DOI 占位。它不是 Nature Reviews Cardiology 的独立摘要目录首页，也不规定其他文体必须采用相同结构。\par
\begin{samepage}在线发表日期与卷期日期可能不同。真实稿件应分别记录来源中的日期字段，页脚使用独立中文译页；不将样板页数或示例日期冒充原刊信息。\par\end{samepage}
\section*{完整覆盖说明占位}
正式交付时，可以在此说明已包含哪些源材料，并对确实缺失的补充文件作出简短界定。覆盖说明应与实际文件一致，不以“全文”一词掩盖未处理的图注或附录。本示例没有原始论文，也没有对任何科学内容作出翻译完成的声明。\par
\SampleBibliography
\SampleSourceInfo
\end{document}
